
\subsubsection{Research Questions}

\textbf{RQ1: Measurement Reliability} - Can Hutchinson trace (10 probes) and power iteration (5 iterations) provide stable rank measurements with coefficient of variation below 15\%?

\textbf{RQ2: Controllability} - Can stable rank be reduced by at least 20\% via gradient-based regularization while maintaining perplexity within 1\% of baseline?

\textbf{RQ3: Isolation} - Does baseline training converge successfully under identical conditions, confirming any failure is specific to the regularization?

\subsubsection{Gate Validation Criteria}

Success required all four criteria:
\begin{enumerate}
\item Mean stable rank reduction ≥20\% vs baseline
\item Perplexity deviation ≤1\% from baseline
\item Layer variance <2.0× mean stable rank
\item Measurement CV <15\%
\end{enumerate}

Failure on any criterion constitutes hypothesis refutation.

\subsubsection{Evaluation Metrics}

\begin{itemize}
\item \textbf{Perplexity}: Evaluated on C4 validation set using sliding window (stride 256)
\item \textbf{Stable Rank}: Measured per layer every 1000 steps via same Hutchinson + power iteration (10 probes, 5 iterations)
\item \textbf{Layer Variance}: Coefficient of variation of stable ranks across 12 layers
\item \textbf{Measurement CV}: Coefficient of variation across Hutchinson samples
\end{itemize}
